\input python_header

\begin{document}
%%%%%%%%%%%%%%%%%%%%%%%%%%%%%%%%%%%%%%%%%%%%%%%
\docheader{2026}{Python}{\S 4 Typed}{Martin Henz, Boyd Anderson}
%%%%%%%%%%%%%%%%%%%%%%%%%%%%%%%%%%%%%%%%%%%%%%%

\input python_intro

\section{Changes}

Python \S 4 Typed is Python \S 4 with type annotations, which are checked before the
program runs. Compared to Python \S 4, Python \S 4 Typed has the following changes:
\begin{itemize}
  \item Function definitions annotate the types of their parameters and their return type,
    and assignments can declare the type of a name.
  \item Function definitions can have type parameters, and the new \lstinline{type}
    statement defines type aliases.
  \item Programs are checked for type errors before they run, as described in Section
    ``Static Type Checking'' below. A program with a type error does not run.
\end{itemize}
A Python \S 4 Typed program without its annotations, type parameters and \lstinline{type}
statements is a Python \S 4 program, and it behaves in the same way.
\input python_4_typed_bnf

\newpage

\input python_indentation
\input python_precedence_3
\input python_restrictions_3


\input python_loops

\input python_names_lang
\input python_numbers

\input python_strings

\input python_static_typing

\input python_typing_front
\input python_typing_middle_34
\input python_typing_back
\input python_typing_tail_34

\input python_standard

\input python_typed_pointer

\input python_misc

\input python_math

\input python_linked_list_library_typed

\input python_pair_mutators

\input python_stream_library_typed

\input python_interpreter

\input generated/python_typed_names

\input python_py_differences

\end{document}
